% !TEX root = ../IPLeiriaMain.tex

% On utilise l'environnement "center" pour contenir l'effet de centrage
% et s'assurer qu'il n'affecte pas le reste du document.
\begin{center}
    \begin{longtable}{|p{3.5cm}|p{11.5cm}|}
        
        % Ligne horizontale en haut du tableau
        \hline
        
        % En-têtes des colonnes en gras
        \textbf{Abréviation} & \textbf{Description} \\
        
        % Ligne horizontale sous les en-têtes
        \hline
        \endfirsthead % Fin de l'en-tête pour la première page
        
        % En-tête pour les pages suivantes (si le tableau continue)
        \hline
        \textbf{Abréviation} & \textbf{Description} \\
        \hline
        \endhead
        
        % Ligne de fin pour chaque page
        \hline
        \endfoot
        
        % Ligne finale en bas du tableau
        \hline
        \endlastfoot
    
        % Contenu du tableau, avec une ligne horizontale après chaque entrée
        API & Application Programming Interface (Interface de Programmation d'Application). \\
        \hline
        PBADS & Personal Behavior Anomaly Detection System (Système de Détection d'Anomalies Comportementales Personnelles). \\
        \hline
        ML & Machine Learning (Apprentissage Automatique). \\
        \hline
        AI & Artificial Intelligence (Intelligence Artificielle). \\
        \hline
        JWT & JSON Web Token (Jeton d'authentification basé sur JSON). \\
        \hline
        CRUD & Create, Read, Update, Delete (Créer, Lire, Mettre à jour, Supprimer). \\
        \hline
        DB & Database (Base de Données). \\
        \hline
        EMSI & École Marocaine des Sciences de l'Ingénieur. \\
        \hline
        JPA & Java Persistence API (utilisé via Spring Data JPA pour l'accès et la gestion des données). \\
        \hline
        REST & Representational State Transfer, style architectural pour les APIs exposées par le backend. \\
        \hline
        PostgreSQL & Système de gestion de base de données relationnelle utilisé pour stocker les données. \\
        \hline
        Spring Boot & Framework Java pour le développement backend de l'application. \\
        \hline
        Spring Cloud & Suite d'outils pour le développement d'applications microservices. \\
        \hline
        Eureka & Service de découverte et registre des microservices (Spring Cloud Netflix Eureka). \\
        \hline
        Kafka & Apache Kafka, broker de messagerie distribué pour la communication asynchrone entre services. \\
        \hline
        Thymeleaf & Moteur de template côté serveur pour le rendu des pages HTML. \\
        \hline
        Chart.js & Bibliothèque JavaScript pour la création de graphiques interactifs. \\
        \hline
        UI / UX & User Interface / User Experience (Interface Utilisateur / Expérience Utilisateur). \\
        \hline
        UML & Unified Modeling Language (Langage de Modélisation Unifié), utilisé pour la modélisation du système. \\
        \hline
        Isolation Forest & Algorithme d'apprentissage automatique non supervisé pour la détection d'anomalies. \\
        \hline
        One-Class SVM & Support Vector Machine à une classe, algorithme alternatif pour la détection d'anomalies. \\
        \hline
        MLflow & Plateforme open-source pour la gestion du cycle de vie du machine learning. \\
        \hline
        scikit-learn & Bibliothèque Python pour l'apprentissage automatique. \\
        \hline
        Gemini AI & Modèle d'intelligence artificielle de Google utilisé pour la génération de recommandations. \\
        \hline
        Resilience4j & Bibliothèque Java pour la tolérance aux pannes (circuit breaker, retry, fallback). \\
        \hline
        BCrypt & Algorithme de hachage de mot de passe utilisé pour la sécurité. \\
        \hline
        Redis & Système de stockage clé-valeur en mémoire utilisé pour le cache et la gestion de session. \\
        \hline
        Prometheus & Système de surveillance et d'alerte open-source pour les métriques. \\
        \hline
        Grafana & Plateforme open-source de visualisation et d'analyse de métriques. \\
        \hline
        HTTP & HyperText Transfer Protocol (Protocole de Transfert HyperTexte). \\
        \hline
        JSON & JavaScript Object Notation (Format de données structuré). \\
        \hline
        CSV & Comma-Separated Values (Valeurs séparées par des virgules). \\
        \hline
        Maven & Outil de gestion et d'automatisation de projets Java. \\
        \hline
        dotenv-java & Bibliothèque Java pour le chargement de variables d'environnement depuis des fichiers .env. \\
        \hline
        WebFlux & Framework réactif de Spring pour les applications non-bloquantes. \\
        \hline
        OpenFeign & Client HTTP déclaratif pour la communication entre microservices (Spring Cloud OpenFeign). \\
        \hline
    
    \end{longtable}
\end{center}
